\section[\appendixname~\thesection]{What Each Arm's Timed Lookup Includes}\label{app:lookup}

For every arm, the timed loop calls the arm's own pass if it has one, as the Go arms do. Otherwise it
runs the harness's loop over the compact copy of the ring. The loop adds each route number, and each
captured value's length and first byte, to a sink that the compiler cannot remove. For an arm without
a 405 of its own, a ``no route'' answer starts a lookup of the same path under every other method
that has routes. Any hit makes the answer 405. On the all-hit rings of the grid, that fallback runs
only after a wrong answer, but its test runs on every lookup. The allocation pass and the latency
samples call each lookup through the harness, so a Go arm's latency includes one call across cgo.

``Index'' in Table~\ref{tab:lookup} means that the harness passes the method as a number and the
adapter indexes one router per method. ``String'' means that the router takes the method's name.

\begin{table}[htbp]
\caption{The timed call of each arm. ``Loop'' is the loop that times the arm, and ``null'' the null
that H2 subtracts.\label{tab:lookup}}
\begin{adjustwidth}{-\extralength}{0cm}
\centering\scriptsize\setlength{\tabcolsep}{3pt}
\begin{tabularx}{\fulllength}{@{}l X l X l l@{}}
\toprule
\textbf{Arm} & \textbf{Timed call} & \textbf{Method} & \textbf{Captured values} & \textbf{405} &
\textbf{Loop; null} \\
\midrule
\arm{regexmatcher-v2} & \FindVtwo, out of line, which calls \texttt{matcher::route::find} & index & views into the path & its own & harness; \arm{null} \\
\arm{regexmatcher-v2-ct} & the same, on the compile-time table & index & views & its own & harness; \arm{null} \\
\AblationArm & the pre-release's lookup, out of line & index & views & its own & harness; \arm{null} \\
\arm{regexmatcher-v1} & the wrapper's match, which calls v1's matching with groups & index & offsets; a new result vector per call & harness & harness; \NullNoFive \\
radix, hash, regex & the harness's lookups & index & views (hash: none) & harness & harness; \NullNoFive \\
Crow & \texttt{crow::Trie::find} & index & copied into strings; views of them & harness & harness; \NullNoFive \\
Oat++ & \texttt{Router::getRoute} & index & views; the tail copied into a string & harness & harness; \NullNoFive \\
Boost.URL & \texttt{parse\_path}, then the router's \texttt{find} & index & views & harness & harness; \NullNoFive \\
r3 & one match entry created, matched and freed per lookup & a bit in the entry & views & harness & harness; \NullNoFive \\
Drogon & the transcribed lookup: a lower-cased copy, two maps, then a regular expression per route & index & a new vector of strings per match & harness & harness; \NullNoFive \\
Pistache & \texttt{sanitizeResource}, then \texttt{findRoute} & index & copies, its only public read & harness & harness; \NullNoFive \\
matchit & \texttt{Router::at} behind one C call & index & offsets & harness & harness; \NullNoFive \\
actix-router & \texttt{Path::set}, then \texttt{Router::recognize}, behind one C call & index & offsets, not decoded & harness & harness; \NullNoFive \\
path-tree & \texttt{PathTree::find} behind one C call & index & offsets & harness & harness; \NullNoFive \\
uWebSockets & \texttt{HttpRouter::route}; the handler stores the route & string & views; none for a catch-all & harness & harness; \NullNoFive \\
glaze & \texttt{route\_table::match} & a switch & decoded strings in a map; one find per name & harness & harness; \NullNoFive \\
cpp-httplib & the request's path assigned, then the method's matchers in order & index & named values copied; groups as views & harness & harness; \NullNoFive \\
httprouter & \texttt{Router.Lookup}, then the handler, in Go & string & offsets & the harness's rule, in Go & its own; its own \\
gin & the vendored tree's search, then the handlers, in Go & string & offsets; values after a return copied & the harness's rule, in Go & its own; its own \\
net/http & \texttt{ServeMux.ServeHTTP} on one reused request, in Go & string & the handler copies each value & its own & its own; its own \\
chi & the context reset, the method map, the tree search, then the handler, in Go & string & offsets & the harness's rule, in Go & its own; its own \\
\bottomrule
\end{tabularx}
\end{adjustwidth}
\end{table}

What remains unequal is stated rather than removed, because removing it would leave the router's
interface:
\begin{itemize}
\item The copies of captured values that a router's interface makes are in its time. This holds for
Crow, Oat++'s tail, Drogon, Pistache, glaze and cpp-httplib's named values. v2 returns views.
\item The adapters of Oat++ and Pistache index one router per method, where the framework looks the
method up in a map. This removes a step from those arms and never adds one. Boost.URL's example
router, matchit and path-tree have no methods, so their adapters keep one router per method.
\item gin's adapter copies the values of a query that needed a return. Only \shape{mixed-overlap} has
such queries, and there gin is incomparable.
\item v2's lookup is an out-of-line call, while a header-only competitor can be inlined into the
timed loop. This adds work to v2 only.
\item The Go arms take the method as a string, inside the timed lookup.
\item r3 creates and frees one match entry per lookup, as its interface requires.
\end{itemize}

\section[\appendixname~\thesection]{Competitor Configurations}\label{app:config}

The C and C++ arms are built by clang \ClangVersion{} in CMake's Release configuration, without
link-time optimization. Every target gets the same flags, \FlagsCxx. The libraries that the harness fetches
and builds get the same flags. The Rust arms use \FlagRust{} and the harness's release profile, with
the prebuilt standard library. The Go arms use \FlagGo, the highest level that L supports, and the C
that cgo compiles gets the native flag too. Table~\ref{tab:config} gives what each competitor adds and
the documentation it follows. ``Library default'' means that we found no documented advice.

\begin{table}[htbp]
\caption{Each competitor's configuration and its source.\label{tab:config}}
\begin{adjustwidth}{-\extralength}{0cm}
\centering\scriptsize\setlength{\tabcolsep}{3pt}
\begin{tabularx}{\fulllength}{@{}l X X@{}}
\toprule
\textbf{Competitor} & \textbf{Build and run-time configuration} & \textbf{The library's documentation} \\
\midrule
Crow \VerCrow & header only, standalone Asio; log level Warning & no optimization advice; its setup page compiles without flags; its logging guide sets Warning \\
Oat++ \VerOatpp & its own CMake, static; its object counters disabled & its build file says the option disables object counting for release builds, for performance \\
Boost.URL \VerBoostUrl & the library and its example router, static & library default \\
r3, \VerRThree & its sources compiled by the harness; PCRE2 \VerPcre{} by its CMake, static, without JIT & no build flags in its README; r3 never compiles a pattern with PCRE2's JIT, so the JIT would not apply \\
Drogon \VerDrogon & its lookup transcribed into the adapter, common flags; the file's hash and line ranges pinned & its guide asks for a Release build of the library, which the harness does not build \\
Pistache \VerPistache & its own CMake, Release, static, tests and TLS off & its README uses meson's release build; CMake's Release is used here \\
uWebSockets \VerUWebSockets & header only; only the router compiled; native flags & no flags in its README; its own example build uses a native build, \FlagOthree{} and link-time optimization \\
glaze \VerGlaze & header only; its CMake adds nothing else on Linux; native flags & its performance guide, written for its serialization, recommends \FlagOtwo{} or \FlagOthree{} and a native build when the program runs where it was built \\
cpp-httplib \VerCppHttplib & header only, no TLS, no compression & library default \\
matchit \VerMatchit & common Rust & library default \\
actix-router \VerActixRouter & common Rust, default features & library default \\
path-tree \VerPathTree & common Rust & library default \\
httprouter \VerHttprouter & common Go & library default \\
gin \VerGin & common Go; the tree vendored and called without gin's engine & its deployment page sets the release mode, which governs debug output only; the vendored lookup reads no mode \\
net/http, Go \VerNetHttp & common Go & no build advice \\
chi \VerChi & common Go; its sources vendored, each with its hash & library default \\
\bottomrule
\end{tabularx}
\end{adjustwidth}
\end{table}

Link-time optimization stays off for every C and C++ arm, v2 included. Of the documentation read, only
uWebSockets' example build uses it.

\section[\appendixname~\thesection]{The Ablation Arm per Cell}\label{app:ablation}

\begin{table}[htbp]
\caption{v2 against the ablation arm, which keeps the pre-release's build and layout. Each value is the
median over the \RoundTwoReps{} pairs of v2's value over the ablation arm's. The run decides nothing
from these ratios.\label{tab:ablation}}
\centering\footnotesize
\begin{tabular}{@{}llrrrr@{}}
\toprule
\textbf{Shape} & $m$ & \textbf{Time} & \textbf{Net instructions} & \textbf{Build time} & \textbf{Heap bytes} \\
\midrule
\AbRowStaticTen \\ \AbRowStaticHundred \\ \AbRowStaticThousand \\ \AbRowStaticTenThousand \\
\AbRowStaticHundredThousand \\ \addlinespace
\AbRowParamLastTen \\ \AbRowParamLastHundred \\ \AbRowParamLastThousand \\ \AbRowParamLastTenThousand \\
\AbRowParamLastHundredThousand \\ \addlinespace
\AbRowParamFirstTen \\ \AbRowParamFirstHundred \\ \AbRowParamFirstThousand \\
\AbRowParamFirstTenThousand \\ \AbRowParamFirstHundredThousand \\ \addlinespace
\AbRowRestTen \\ \AbRowRestHundred \\ \AbRowRestThousand \\ \AbRowRestTenThousand \\
\AbRowRestHundredThousand \\ \addlinespace
\AbRowWildTen \\ \AbRowWildHundred \\ \AbRowWildThousand \\ \AbRowWildTenThousand \\
\AbRowWildHundredThousand \\ \addlinespace
\AbRowMixedDisjointTen \\ \AbRowMixedDisjointHundred \\ \AbRowMixedDisjointThousand \\
\AbRowMixedDisjointTenThousand \\ \AbRowMixedDisjointHundredThousand \\ \addlinespace
\AbRowMixedOverlapTen \\ \AbRowMixedOverlapHundred \\ \AbRowMixedOverlapThousand \\
\AbRowMixedOverlapTenThousand \\ \AbRowMixedOverlapHundredThousand \\
\bottomrule
\end{tabular}
\end{table}
